\scp{Unclassified languages}{07-05}
Before we discuss each of the subgroups within central
Maluku in more detail,
we need to draw attention to two languages which
cannot yet be classified on the basis of shared phonological
innovations due to lack of data.
These are \ili{Salas}, of eastern Seram and \ili{Hukumina} of northwest \ili{Buru}.

\sscp{\ili{Salas} [sgu]}{07-05-01}
According to \citet[121f]{LoskiLoski1989} “\ili{Salas} Gunung is
spoken by 50 people in the village by that name in the \ili{Waru} Bay.
Most of the people are reported to be bilingual in the \ili{Masiwang} language.
\citet{Tauern1931} presents a short grammar of \ili{Liambata},
possibly a dialect of \ili{Salas} Gunung”.
\ili{Salas} is placed by \citet{LoskiLoski1989} within their “\ili{Manusela}-\ili{Seti}
language chain” and thus within our \tsc{\ili{Ambon}-Seram},
and not with \ili{Bobot} or \ili{Masiwang},
which belong in our \tsc{Seram-Tanimbar-Bomberai}.

Our examination of the ``\ili{Liambata}'' data in \citet[128--139]{Tauern1931}
indicates that it may be a variety of \ili{Seti},
based on distinctive lexical items such as \it{inta} ‘skin’ (\ili{Liana-Seti}
\it{inta} ‘skin’), and \it{ufeia} ‘house’ (\ili{Liana-Seti} \it{ufai-a} ‘house’).
At the same time,
this data diverges from other sources of \ili{Seti} in several respects
indicating that it may indeed be a different language.
It is important to note that the data in \citet[128--139]{Tauern1931}
is different from the “\ili{Liambata}” data in \citet{Stresemann1927},
the latter of which \it{is} (almost certainly) a variety of \ili{Seti}.

Given the equivocation of \citet{LoskiLoski1989} regarding the
data in \citet{Tauern1931}, as well as our own judgment that the data from \citet{Tauern1931}
is related to but different from \ili{Seti},
we cannot confidently identify any \ili{Salas} data to make an assessment of its position.
Lexically, it is closer to \tsc{\ili{Ambon}-Seram} languages than \tsc{Seram-Tanimbar-Bomberai}.

\sscp{\ili{Hukumina} [huw]}{07-05-02}
\ili{Hukumina} is an extinct language of north-western \ili{Buru}.
\citet[63--76]{Stokhof1982b} contains a word list labelled ``\ili{Hukumina}'',
but the reliability of this list is very low,
and it contains many items which are clearly \ili{Buru} and \ili{Kayeli}.
Grimes was able to collect a small amount of data in 1989
and states:

\begin{quote}
	``\ili{Hukumina} was referred to by \citet{Hendriks1897} as a dialect of \ili{Buru} spoken in the districts
	of \ili{Hukumina}, Palumata and Tomahu (these were districts in the north-west comer of \ili{Buru}).
	He gives no other details other than his impression that one could shift easily between
	\ili{Hukumina} and other dialects of \ili{Buru}. A \ili{Hukumina} word list is published in the Holle lists
	(Stokhof 1982) [\cite{Stokhof1982b}] with very little information other than the date 1896.
	
	[\ldots] I was able to find one old woman from the former village of \ili{Hukumina} that used to be
	located behind the present village of \ili{Masarete} near the fort at \ili{Kayeli} (north-\ili{east} \ili{Buru}). My
	data are even more divergent than the Holle list (the latter includes several items that are
	clear borrowings from \ili{Buru}). The language of \ili{Hukumina} is very distinct from \ili{Buru} in
	intonation, lexicon and sound correspondences. For example, PAn *k is \ili{Buru} /k/ and
	\ili{Hukumina} /c/ intervocalically.'' \citep[99]{Grimes2000a}
\end{quote}

\begin{table}
	\centering
	\caption[\ili{Buru}, \ili{Kayeli}, and \ili{Hukumina} words (Grimes 2010: 78)]
	{\ili{Buru}, \ili{Kayeli}, and \ili{Hukumina} words \citep[78]{Grimes2010b}}
	\label{tab:07-19}
	\st{6pt}\begin{tabular}{llll}\lsptoprule
Gloss	&	\ili{Buru}	&	\ili{Kayeli}	&	\ili{Hukumina}	\\	\midrule
hand, arm	&	fahan	&	limani	&	rima	\\	
mouth	&	fifin	&	nuan	&	mpidu-	\\	
foot, leg	&	kadan	&	bitiɲ	&	eɲaik	\\	
head	&	olon	&	oloni	&	fatun	\\	
betelpepper	&	dalu	&	gamal	&	elut	\\	
jungle	&	mua	&	boot	&	gulalin	\\	
sleep	&	bage	&	ine	&	etnono	\\	
sleepy	&	duba	&		&	lastoton	\\	
sick, painful	&	empei	&	berere	&	boto	\\	
return (home)	&	oli	&	wae	&	bihin	\\	
no, not	&	moo	&		&	mbaa	\\	
hungry	&	eglaba	&	nanibar	&	spala	\\	
sit	&	defo	&	stea	&	egnabon	\\	
		\lspbottomrule
	\end{tabular}
\end{table}

He further states that the consultant's ``[\ldots] mind wandered
regularly, and the little data I was able to collect from her are a mixture of \ili{Kayeli},
\ili{Buru}, and some other language that I assume is \ili{Hukumina}.'' \citep[75]{Grimes2010}
A sample of \ili{Hukumina} vocabulary compared with \ili{Buru} and \ili{Kayeli} is given in \trf{07-19}
to illustrate the difference between this language and others of \ili{Buru} Island.

Grimes' suspicion is that \ili{Hukumina} may have been a long-standing colony from
South\ili{east} Sulawesi, as there are thousands of migrants from there along the
west and south coasts of \ili{Buru}. However, the lack of reliable data means
that this cannot be confirmed. Unless more data becomes available,
\ili{Hukumina} remains unclassifiable. We list it with other \tsc{Sula-Buru} languages
in \crf{ch:SB} and appendix \ref{ch:LanLis} out of convenience according to its geographic location.
